\documentclass[12pt, letterpaper]{article}
\usepackage{graphicx} % Required for inserting images
%\usepackage[T1]{fontenc}

\usepackage{amsfonts}
\usepackage{tabularx}
\usepackage{multirow}
\usepackage{booktabs}
\usepackage{float}
\usepackage{caption}
%\usepackage{natbib}
\title{Programy użytkowe \\ Lista 7}
\date{}
\usepackage{polski}
\begin{document}
\maketitle



\begin{enumerate}
    \item Korzystając z \texttt{set polar} narysuj spiralę Archimedesa. Zwróć uwagę na nową nazwę zmiennej. Ustaw jej zakres tak, aby nanieść 3 ``zwoje'' spirali na wykres.
    \item Narysuj okrąg o promieniu zapisanym we wcześniej utworzonej zmiennej. Wykorzystaj współrzędne biegunowe (\texttt{set polar}) lub krzywą zadaną parametrycznie (\texttt{set parametric}).
\item Zapisz zadanie 1. w formie skryptu. Nich skrypt przyjmuje argument określający długości półosi wielkiej oraz małej oraz zapisuje wykres w pliku. Wskazówka: \texttt{gnuplot -c} oraz \texttt{argv[]}.
\item Zapisz zadanie 2. w formie skryptu. Niech skrypt przyjmuje jako argument liczbę ``zwojów''x spirali do narysowania.
\end{enumerate}

\textbf{Zadanie domowe.} Wyniki zadań 1. oraz 2. załącz w notatce w \LaTeX{-u}, zapisując również wzory opisujące otrzymane krzywe.


\end{document}